\section{Discusión}
\label{sec:discusion}

\paragraph{Qué compra realmente cada mecanismo.}
El recorte de leyenda y el rechazo por confianza cambian los dos cobertura por calidad macro, pero
le cobran a clientes distintos. Recortar la leyenda rechaza un tipo de cultivo entero en todas
partes: quien lo siembra no recibe nada, en ninguna parcela, en ninguna región. La abstención
rechaza parcelas sueltas repartidas entre todos los cultivos, así que ningún productor queda
excluido de forma sistemática, que es la propia preocupación de la literatura de clasificación
selectiva sobre quién paga la abstención \cite{jones2020selectivedisparities}. Nuestras mediciones
dicen además que la abstención es la mejor compra a igual cobertura. Los dos argumentos apuntan en
la misma dirección, y no conocemos escenario operativo en que sea preferible lo contrario salvo
uno: cuando el consumidor aguas abajo exige un catálogo fijo y publicado y no puede actuar sobre un
mapa parcial.

\paragraph{Qué habría que reportar.}
Una cifra promediada por clase calculada sobre un catálogo podado es ininterpretable por sí sola, y
el arreglo no es caro ni novedoso. Repórtese el número de clases sobre el que se toma el promedio,
la cobertura que el producto entrega, y el control de la sección~\ref{sec:control}. Tres números,
uno de los cuales cuesta una sola línea de código. Sin ellos, un artículo que sube el F1 macro un
cuarto de punto tras retirar tres clases raras está reportando aritmética, y quien lo lee no puede
distinguirlo de un artículo que mejoró un modelo.

\paragraph{Sobre el criterio que se usa en la práctica.}
Que la retirada por soporte sea a la vez el criterio más débil y no monótona merece énfasis, porque
el soporte es el criterio que se usa. Se elige por una razón defendible —una clase con pocas
muestras es una clase cuya estimación es inestable— pero inestabilidad de la estimación y mala
separabilidad son propiedades distintas, y el soporte solo rastrea la primera. Una clase con
\num{1718} parcelas que el modelo separa limpiamente aporta más a un mapa entregado que dos clases
abundantes que el modelo no logra distinguir. Elegir por calidad por clase medida, sobre bloques
disjuntos de aquellos donde se evaluará, no cuesta nada más y es estrictamente mejor en todas las
comparaciones que corrimos.

\paragraph{Sobre cómo medir meta-modelos.}
Las tres cifras del mismo modelo de apilamiento —\num{0.7470}, \num{0.6789} y \num{0.5360}— no son
tres estimaciones plausibles entre las que elegir. La primera está contaminada, y la formulación
original del apilamiento ya exigía las particiones libres de fuga que ella viola
\cite{wolpert1992stacking}. La tercera está libre de fuga pero castiga al modelo por un artefacto
de agregación: promediar macros por bloque en un escenario desbalanceado puntúa cada bloque sobre
clases que sus parcelas no contienen. Agrupar las posteriores fuera de fold y puntuar una sola vez,
que da \num{0.6789}, es el estimador que responde a la pregunta que de verdad se está haciendo. La
distancia entre la primera y la última es lo bastante ancha como para invertir una conclusión sobre
si el ensamble aporta algo, y por eso el régimen tiene que declararse cada vez que se imprime una
cifra así.
